% _preamble_shared.tex —— 单章(_tmp.tex) 与整书(main.tex) 模板共用的「补丁片段」唯一真源。
%
% 这是消除「双源漂移」的关键：过去中文 URL 支持、参考文献降级、codeblock 环境、
% 防溢出、表格字号这 5 段补丁在 _tmp.tex 与 main.tex 里各写一份，任何人改一处都极易漏另一处。
% 现在它们只在此文件定义一份；fix / scaffold 在生成 _tmp.tex / main.tex 时把
% 模板里的共用补丁占位符替换为本文件内容（内联，保证编译产物自包含、无 \input 路径问题）。
%
% 改这些补丁，只改这里，重新 fix / scaffold 即可，两套模板自动同步。
% 注意：本片段只包含「补丁命令」，不含 \usepackage（包清单仍由各模板各自声明并保持同步）。

% === 中文 URL 支持 ===
% url.sty 默认在 math mode 渲染 URL，xeCJK 不接管 math mode，中文会变成豆腐块。
% 重写 \Url@FormatString 去掉 math mode，中文走正常字体分类即可正确显示。
\makeatletter
\def\Url@FormatString{%
 \UrlFont
 \expandafter\UrlLeft\Url@String\UrlRight
}
\makeatother

% === 参考文献降级 + 不跳页 + 不改页眉 ===
% book/ctexrep 的 thebibliography 默认 \chapter*（最大级标题 + 跳页 + 改页眉）。
% 重定义为 \section* 并去掉 \@mkboth，与正文页眉保持一致。
\makeatletter
\patchcmd{\thebibliography}{\chapter*{\bibname}}{\section*{\bibname}}{}{}
\patchcmd{\thebibliography}{\@mkboth{\MakeUppercase\bibname}{\MakeUppercase\bibname}}{}{}{}
\makeatother

% === 代码块环境：浅灰打底 + 自动换行（禁用裸 \verbatim） ===
% tcolorbox + listings 引擎（\newtcblisting），breaklines 自动折行长代码行，
% 与正文明显区分。禁用裸 \verbatim（不折行会溢出页面）。
% 代码字号比正文小一档（\footnotesize）：长代码行更不易溢出，视觉上也与正文拉开层级。
\newtcblisting{codeblock}{
  colback=gray!10!white, colframe=gray!45!white, boxrule=0.4pt,
  arc=3pt, left=6pt, right=6pt, top=5pt, bottom=5pt,
  listing only, breakable,
  listing options={
    breaklines=true, breakatwhitespace=false,
    basicstyle=\footnotesize\ttfamily, showstringspaces=false,
    frame=none, numbers=none, aboveskip=0pt, belowskip=0pt
  }
}

% === 表格字号：略小于正文 ===
% 三线表单元格内容统一 \small（11pt 正文 → 10pt 表格）。
\AtBeginEnvironment{tabularx}{\small}
\AtBeginEnvironment{tabular}{\small}
\AtBeginEnvironment{longtable}{\small}

% === 防溢出 ===
\setlength{\emergencystretch}{3.5em}
